\section{Production-Grade Implementation}
\label{sec:implementation}

% THIS IS THE SI-FIT SECTION. The MDPI Applied Sciences "Innovations in Supply
% Chain Resilience" call wants digital-technology contributions to risk
% mitigation. This section argues that production-ready, hardened, open-source
% optimisation tooling IS the digital-innovation contribution — not a
% deployment afterthought.
%
% Outline (target ~2500-3500 words; this is one of the two largest sections):
%
% 4.1 Design goals
%   - Reproducibility, auditability, security posture, deployability.
%   - Why these matter for SC resilience adoption (vs research-grade scripts).
%
% 4.2 Architecture
%   - Package layout: src/presidio_vol_assign/ (CLI entry point `pva`), with a
%     pluggable Domain adapter per problem family in domains/: ed_staffing.py
%     (v0.1.0, Rabiei et al. ESWA 2023) and humanitarian.py (v0.2.0, this paper).
%   - FIS engine (fis.py, fis_humanitarian.py), shared EA engine (engine.py),
%     CLI surface (cli.py).
%   - Pluggable index/objective interfaces via the Domain ABC (domains/base.py).
%
% 4.3 Engineering controls
%   - Security: PRESIDIO-REQ.md threat model; CodeQL v4 in CI; SECURITY.md
%     vulnerability disclosure; unified SDLC pointer.
%   - Reproducibility: uv.lock for deterministic Python dependency resolution;
%     fixed Python floor (>=3.10); deterministic RNG seeding for EAs.
%   - Distribution: PyPI package, semantic versioning (v0.1.0 -> v0.2.0
%     reflects the addition of affected-people allocation), OSS license.
%   - Documentation: README, instruction PRESIDIO-REQ.md, type hints, doctests.
%
% 4.4 The v0.2.0 extension (this paper's tooling delta)
%   - New modules: domains/humanitarian.py + fis_humanitarian.py.
%   - Three new FIS implementations (FLPP, TFL, CABL).
%   - New CLI subcommand: `pva allocate-people` (alias for
%     `pva assign --model humanitarian`); `pva benchmark`, `pva sensitivity`,
%     and `pva show` for evaluation/figures.
%   - New test fixtures: small (5 RCs / 150 people) and large (10 RCs / 300
%     people) cases (generated by `pva benchmark`; see examples/).
%
% 4.5 Reproducibility contract
%   - Reproduction is deterministic via seeded RNG from CSV inputs (no YAML
%     --config). The committed examples/ datasets (small = 12/3, paper_scale =
%     150/5; regenerate via examples/generate_examples.py) plus `pva benchmark`
%     regenerate every table; `pva show` regenerates the figures. E.g.:
%       pva benchmark --model humanitarian --size both --seed 42 --check-repro
%       pva allocate-people --people <people.csv> --centers <centers.csv> \
%                           --solver nrga-ranked --seed 42 --output results/

TODO.

\subsection{Command-line interface}
\label{sec:cli}

The affected-people allocation model is exposed as the \texttt{pva
allocate-people} subcommand (a convenience alias for \texttt{pva assign
--model humanitarian}). A complete run, including the canonical NRGA solver of
Section~\ref{sec:solvers}, is invoked as:
%
\begin{quote}\small\ttfamily
pva allocate-people --people people.csv --centers centers.csv \\
\hspace*{2em}--solver nrga-ranked --seed 42 --output results/
\end{quote}
%
The \texttt{--solver} flag accepts \texttt{nsga2}, \texttt{nrga} (lightweight
variant), \texttt{nrga-ranked} (canonical, rank-biased roulette; used for the
reported NRGA results), \texttt{both} (\texttt{nsga2}+\texttt{nrga}), and
\texttt{all}. The \texttt{--seed} flag fixes the RNG for the bit-for-bit
reproducibility evaluation of Section~\ref{sec:evaluation}.

Two further subcommands support the evaluation: \texttt{pva benchmark}
regenerates the small/large instances and emits the Table-3-style mean$\pm$std
summary (with an optional bit-for-bit \texttt{--check-repro} column), and
\texttt{pva sensitivity} sweeps multiplicative perturbations of the FIS output
scores (e.g.\ $\pm 10\,\%$, $\pm 20\,\%$) to report how the Pareto metrics shift
under rule-base uncertainty (Section~\ref{sec:evaluation}).
